\documentclass[12pt]{article}

\usepackage{sbc-template}
\usepackage{graphicx,url}
\usepackage[utf8]{inputenc}
\usepackage[brazil]{babel}
\usepackage{float}

\sloppy

\title{DLM-PDF: Uma Plataforma Web para a Transposição da Posse Física para o Livro Digital via Blockchain}
\author{Arnaldo M. Quintanilha\inst{1}, Gabriel [Sobrenome]\inst{1}, Matheus P. Raposo\inst{2}}

\address{
  $^1$Bacharelado em Sistemas de Informação -- CEFET/RJ \\
  $^2$Mestre em Ciência e Tecnologia de Materiais -- UERJ \\
  Nova Friburgo -- RJ -- Brazil
  \email{arnaldomartinsq@gmail.com, gabriel@email.com, matheusmerlim@hotmail.com}
}

\begin{document}

\maketitle

\begin{abstract}
This paper proposes a digital ownership model for e-books based on Blockchain technology. The model allows users to resell and lend digital books — rights commonly available in physical markets but blocked by traditional DRM systems. Instead of protecting the file itself, as in Oliveira and Lazarin (2020), the proposed approach protects only the access authorization, validated in real time through a smart contract. A test suite of 31 automated tests was executed, covering license registration, minting, access validation, P2P transfer, digital lending, AES-256-CBC encryption, and HMAC integrity verification. The complete authentication handshake averaged 1.26 ms in isolated execution; with local network latency (Ganache/JSON-RPC), the total validation cycle reaches approximately 1.2 seconds per operation. The main limitation is the need for continuous internet connectivity, since the decryption key is never stored locally.
\end{abstract}

\begin{resumo}
Este trabalho propõe um modelo de posse digital para e-books baseado em Blockchain. O modelo permite que o usuário revenda e empreste livros digitais — direitos comuns no mercado físico, mas bloqueados pelos sistemas de DRM tradicionais. Em vez de proteger o arquivo em si, como em Oliveira e Lazarin (2020), a abordagem proposta protege apenas a autorização de acesso, validada em tempo real por um contrato inteligente. Uma suíte de 31 testes automatizados foi executada, cobrindo registro de livro, emissão de licença, validação de acesso, transferência P2P, empréstimo digital, cifragem AES-256-CBC e verificação de integridade HMAC. O handshake de autenticação completo apresentou média de 1,26 ms em execução isolada; com a latência de rede local (Ganache/JSON-RPC), o ciclo total de validação alcança aproximadamente 1,2 segundos por operação. A principal limitação é a necessidade de conectividade contínua à internet, uma vez que a chave de decriptografia nunca é armazenada localmente.
\end{resumo}

\section{Introdução}

No decorrer dos últimos anos, o mercado de conteúdos digitais vem crescendo de forma expressiva. Livros, músicas e filmes, que antes eram comercializados quase exclusivamente em formato físico, passaram a ser distribuídos por plataformas digitais. Esse processo trouxe consigo um problema que ainda não foi resolvido de forma satisfatória: como garantir ao consumidor o mesmo direito de posse que ele teria sobre um objeto físico?

Quando alguém compra um livro impresso, pode revendê-lo, emprestá-lo ou doá-lo. No ambiente digital, essas práticas são amplamente vetadas. Plataformas como Amazon Kindle e Adobe Digital Editions impõem restrições que, na prática, convertem a compra de um e-book numa licença de uso revogável — e não numa aquisição de posse real sobre o bem.

Os sistemas de Gestão de Direitos Digitais (DRM) foram criados afim de proteger os detentores dos direitos autorais contra cópias e redistribuições não autorizadas (OLIVEIRA, 2020). Entretanto, ao centralizar o controle de acesso em servidores proprietários, esses sistemas acabam por restringir liberdades que o consumidor teria naturalmente no mercado físico.

A tecnologia Blockchain surge como uma alternativa para esse cenário. Através de contratos inteligentes\footnote{Programas armazenados na blockchain que executam automaticamente as regras de uma transação, sem necessidade de intermediário.}, é possível criar um registro imutável da titularidade de um ativo digital, tratando-o de maneira análoga a um objeto físico — único e transferível.

O presente trabalho propõe um modelo baseado na arquitetura de Oliveira e Lazarin (2020) — trabalho desenvolvido no próprio ambiente institucional do CEFET/RJ — adaptado para funcionar via navegador web, sem que o usuário precise instalar qualquer software. O ponto central da proposta é deslocar o controle de segurança: em vez de proteger o arquivo em si, protege-se a autorização de acesso, validada em tempo real na rede blockchain.

\section{Revisão Bibliográfica}

A fundamentação teórica deste trabalho foi construída a partir de revisão bibliográfica em bases de dados como IEEE Xplore, ResearchGate e ACM Digital Library. A primeira busca utilizou os descritores "digital rights management" AND ("resell" OR "resale") AND "e-book", retornando 32, 120 e nenhum artigo respectivamente. Posteriormente, aplicou-se a combinação "Blockchain-based" AND "Authentication Service", com 60 resultados no IEEE Xplore, 172 no ResearchGate e 100 na ACM Digital Library. Os trabalhos com maior aderência ao tema foram selecionados como base desta pesquisa.

\subsection*{Gestão de Direitos Digitais}

A expansão da internet possibilitou a transposição de conteúdos culturais para o meio digital. Com isso, surgiu a necessidade de mecanismos capazes de controlar o acesso a esses conteúdos e impedir sua reprodução não autorizada. As DRM foram desenvolvidas com esse objetivo.

Daewon (2015) classifica esses sistemas em "fortes" ou "fracos" conforme o nível de restrição imposto ao usuário. Os mais restritivos oferecem maior proteção ao detentor dos direitos, mas prejudicam a experiência de quem adquire o produto. Os menos restritivos preservam a liberdade do usuário, porém aumentam o risco de pirataria. Nenhum dos dois extremos resolve o problema de forma equilibrada.

Nair et al. (2005) propuseram um modelo voltado a redes P2P\footnote{Redes ponto-a-ponto, nas quais os próprios usuários distribuem o conteúdo entre si, sem servidor central.}, no qual uma licença especial autoriza a revenda do conteúdo adquirido. O modelo estabelece que cada usuário deve obter uma licença válida antes de ter acesso ao material. O estudo aponta dificuldades técnicas relevantes para sua viabilização: garantir a não-duplicidade do ativo digital, assegurar a rastreabilidade das transações e permitir que o usuário acesse o conteúdo em seus diferentes dispositivos.

Beetz (2008) destaca o potencial do ambiente digital como mercado para produtos intelectuais, mas aponta como um dos pontos críticos a dificuldade de replicar digitalmente o direito de revenda — prática corriqueira e consolidada no mercado físico.

\subsection*{Blockchain e Contratos Inteligentes}

A tecnologia Blockchain tornou-se conhecida principalmente por sua aplicação em criptomoedas, como Bitcoin e Ethereum. Caracteriza-se por armazenar dados em uma cadeia linear e cronológica de blocos, formando um registro distribuído e imutável. Esse sistema dispensa a necessidade de intermediários centrais, como bancos ou plataformas proprietárias.

Gaber e Zhang (2010) descrevem os contratos inteligentes como programas armazenados na blockchain que executam automaticamente os termos de uma transação. A execução autônoma das regras estabelecidas garante transparência sem necessidade de intervenção humana.

Seguindo essa proposta, Oliveira (2020) desenvolveu uma arquitetura estruturada em três componentes: o Controlador, responsável pelo armazenamento dos arquivos e geração de chaves criptográficas; o Leitor, software que valida as requisições e decifra o conteúdo para visualização; e a Blockchain, que assegura a integridade e rastreabilidade das transações. Nesse modelo, o usuário acessa apenas a versão renderizada pelo software leitor. O acesso direto ao arquivo binário original fica impedido.

\subsection*{Autenticação Descentralizada via APIs}

Arya et al. (2020) propuseram uma arquitetura de marketplace de APIs baseada em blockchain, na qual a execução dos modelos é distribuída entre múltiplos fornecedores de nuvem. Nesse ecossistema, os contratos inteligentes funcionam como livro-razão imutável, registrando cada etapa da transação e garantindo resiliência contra comportamentos maliciosos de provedores e consumidores.

Wang e Hu (2018) demonstraram a eficiência de blockchains de consórcio para armazenamento de certificados de autorização. A validação ocorre pela comparação do hash do certificado do usuário com o registro na rede — processo simples e auditável.

Lucas et al. (2025) mostram que a integração de APIs RESTful com contratos inteligentes é capaz de reduzir a complexidade de adoção de sistemas descentralizados pelo usuário final, facilitando a transição de arquiteturas centralizadas para descentralizadas.

Kustov et al. (2022) propõem o uso de um Oráculo DVCS\footnote{Data Validation and Certification Server — servidor que atua como ponte entre sistemas externos e a blockchain.} para validar dados e mediar requisições de acesso, preservando a integridade dos direitos autorais. Esse conceito é central na arquitetura proposta no presente trabalho.

\section{Desenvolvimento}

\subsection{Do Arquivo à Autorização}

O modelo de Oliveira (2020) concentra a segurança na custódia do arquivo original. O usuário precisa de um software de leitura específico instalado localmente, e o conteúdo só pode ser acessado por meio desse programa. A abordagem é eficaz do ponto de vista da proteção, mas limita a flexibilidade do usuário e cria uma dependência de software proprietário que não existia no mercado físico.

A proposta deste trabalho segue um caminho diferente. A segurança não reside no arquivo em si, mas na autorização de acesso. O arquivo binário pode circular livremente — sozinho, ele não contém a chave necessária para ser lido. A posse é um estado registrado na blockchain, não uma propriedade do arquivo.

Sendo assim, qualquer editora pode registrar seus livros na rede e qualquer usuário pode validar sua posse diretamente com o contrato inteligente, sem depender de uma plataforma intermediária. A distribuição do conteúdo passa a ser descentralizada.

\subsection{Arquitetura do Sistema}

O sistema é composto por três camadas que se comunicam por meio de APIs RESTful.

A primeira é o \textbf{Smart Contract} (DLMPDFLicense.sol), implantado na rede Ethereum. É o componente central e imutável do sistema. O contrato mantém, para cada exemplar de um livro digital, o registro de qual endereço de carteira detém a posse naquele momento. A transferência de posse revoga automaticamente o acesso do vendedor, sem intervenção de qualquer servidor externo. É a implementação direta da Regra 1 para 1 que fundamenta o modelo.

A segunda é o \textbf{Servidor de Autenticação}, aplicação Node.js que atua como Oráculo no sentido descrito por Kustov et al. (2022) — um intermediário entre o cliente web e a blockchain. Ao receber uma requisição de abertura de documento, o servidor consulta o contrato inteligente e, caso a posse seja confirmada, retorna uma chave de sessão efêmera ao cliente. Vale notar que essa chave tem validade única: ela não é reutilizável entre sessões distintas.

A terceira camada é o \textbf{Cliente Web}, interface executada no próprio navegador do usuário. É responsável por autenticar a identidade via assinatura criptográfica da carteira Ethereum, receber a chave de sessão e realizar a decriptação do arquivo em memória — sem que o conteúdo seja gravado em disco em nenhum momento.

\subsection{Formato do Arquivo .dlm}

O PDF original é encriptado e distribuído no formato proprietário \texttt{.dlm}. A escolha de um formato próprio — em vez de usar PDF criptografado diretamente — decorre de uma limitação prática: o formato PDF não possui campo nativo para armazenar um identificador de licença vinculado a uma rede externa. Sendo assim, optou-se por um container binário simples com a seguinte estrutura:

\begin{verbatim}
[ MAGIC      4B : "DLM\x01"                   ]
[ licenseId  8B : uint64 big-endian            ]
[ IV        16B : vetor de inicialização AES   ]
[ HMAC      32B : integridade HMAC-SHA-256     ]
[ ciphertext NB : conteúdo AES-256-CBC         ]
\end{verbatim}

O conteúdo é cifrado com AES-256-CBC. A chave de cifragem não fica armazenada no arquivo — ela é derivada a partir de um segredo mantido no servidor de autenticação, utilizando o algoritmo HKDF-SHA-256 com o \texttt{licenseId} como parâmetro de entrada. Sendo assim, copiar o arquivo \texttt{.dlm} sem possuir a licença correspondente na blockchain resulta em um dado ilegível.

O campo HMAC garante a integridade do binário. Dashti et al. (2007) identificaram a necessidade desse mecanismo em protocolos de DRM, afim de prevenir ataques de substituição de conteúdo — situação em que um arquivo adulterado é apresentado ao usuário como se fosse o original registrado pela editora.

\subsection{Fluxo de Acesso}

Quando o usuário seleciona um arquivo \texttt{.dlm} no navegador, o sistema começa lendo apenas os primeiros bytes do cabeçalho para extrair o \texttt{licenseId}. O conteúdo cifrado não é tocado nessa etapa — não há sentido em decriptografar algo antes de confirmar que o acesso é válido.

Com o \texttt{licenseId} em mãos, o cliente solicita ao servidor um nonce\footnote{Número aleatório de uso único, utilizado para evitar ataques de repetição de requisições.} associado ao endereço da carteira Ethereum do usuário. O usuário então assina esse nonce com a chave privada da carteira — operação que prova a identidade sem transmitir a chave em si. O servidor verifica a assinatura e emite um token JWT de sessão, seguindo lógica análoga à descrita por Wang e Hu (2018) para autenticação por certificados em blockchain.

A partir daí, o cliente envia a requisição de abertura ao servidor, que chama a função \texttt{validateAccess()} no contrato inteligente. Diferente de uma simples consulta, essa chamada registra o evento na blockchain — o que significa que toda tentativa de acesso fica gravada de forma imutável, seja ela bem-sucedida ou não. Se o acesso é confirmado, o servidor deriva uma chave de sessão via HKDF e a retorna. O cliente usa a Web Crypto API nativa do navegador para decriptografar o conteúdo em memória, e o PDF é renderizado pela biblioteca PDF.js. Ao fechar a aba, o conteúdo é descartado.

\subsection{Revenda e Empréstimo}

Um dos pilares do modelo é o que foi denominado Regra 1 para 1: em qualquer momento, apenas um endereço possui acesso ativo a um dado exemplar. Ao realizar uma revenda, o contrato inteligente executa a transferência de titularidade e distribui automaticamente os royalties ao endereço da editora, conforme o percentual definido no cadastro do livro. O acesso do vendedor é revogado sem necessidade de qualquer ação do servidor.

O empréstimo segue lógica similar. O proprietário pode ceder o acesso temporariamente a outro endereço, por um período de até 30 dias. Durante o empréstimo, o próprio dono perde o acesso — comportamento que replica o que acontece com um livro físico emprestado. Ao término do prazo, o contrato revoga a cessão automaticamente na próxima consulta. Dentre as aplicações práticas desse mecanismo, destaca-se a possibilidade de criar sebos digitais com comportamento análogo ao do mercado físico.

\subsection{Comparativo com Abordagens Existentes}

A diferença entre o DLM-PDF e os sistemas de NFT convencionais merece atenção. Nesses últimos, possuir o token não garante controle sobre o acesso ao arquivo: o PDF costuma ser distribuído por link público, completamente desvinculado da propriedade do token. Qualquer pessoa com o link consegue acessar o conteúdo. No DLM-PDF, ativo e autorização de acesso são inseparáveis — a chave de decriptografia só é liberada após confirmação on-chain da titularidade do exemplar. Isso resolve a brecha apontada por Nair et al. (2005) em redes P2P, onde licença e conteúdo trafegam de forma dissociada.

Em relação ao DRM tradicional, a diferença está na ausência de plataforma central. Não há servidor de licenças proprietário que possa ser descontinuado, hackeado ou que imponha restrições unilaterais ao usuário. A Tabela~\ref{tab:comparativo} sintetiza essa comparação.

\begin{table}[H]
\centering
\caption{Comparativo entre abordagens de DRM para e-books}
\label{tab:comparativo}
\begin{tabular}{|l|c|c|c|}
\hline
\textbf{Recurso} & \textbf{DRM Tradicional} & \textbf{NFT Comum} & \textbf{DLM-PDF} \\
\hline
Interoperabilidade P2P        & Não & Sim & Sim \\
\hline
Independência de plataforma   & Não & Não & Sim \\
\hline
Prevenção de duplicação       & Média & Baixa & Alta \\
\hline
Revenda com royalty           & Não & Parcial & Sim \\
\hline
Empréstimo com revogação real & Não & Não & Sim \\
\hline
\end{tabular}
\end{table}

\subsection{Resultados dos Testes}

A avaliação do modelo foi conduzida por meio de uma suíte de 31 testes automatizados, implementada em Node.js v22 com simulação completa da lógica do contrato inteligente. Os testes cobriram nove categorias: registro de livro, emissão de licença, oráculo de autenticação, transferência P2P, empréstimo digital, devolução de empréstimo, cifragem AES-256-CBC, handshake completo e carga com 100 validações consecutivas. Todos os 31 casos foram aprovados.

A Tabela~\ref{tab:tempos} consolida os tempos de execução medidos para cada operação principal.

\begin{table}[H]
\centering
\caption{Tempos de execução medidos — execução isolada, sem latência de rede}
\label{tab:tempos}
\begin{tabular}{|l|r|}
\hline
\textbf{Operação} & \textbf{Tempo (ms)} \\
\hline
\texttt{registerBook} & 0,442 \\
\hline
\texttt{mintLicense} & 0,076 \\
\hline
\texttt{validateAccess} — proprietário & 0,048 \\
\hline
\texttt{validateAccess} — acesso negado & 0,017 \\
\hline
\texttt{checkAccess} view (sem gas) & 0,036 \\
\hline
\texttt{transferLicense} (revenda P2P) & 0,040 \\
\hline
\texttt{lendLicense} (empréstimo) & 0,051 \\
\hline
\texttt{returnLicense} (devolução) & 0,037 \\
\hline
Cifragem AES-256-CBC — PDF de 50 KB & 2,014 \\
\hline
Decifragem + verificação HMAC & 0,828 \\
\hline
\textbf{Handshake completo} (nonce → JWT → sessionKey) & \textbf{1,260} \\
\hline
100 validações consecutivas (carga) & 0,133 (total) \\
\hline
\end{tabular}
\end{table}

O handshake completo — ciclo que engloba a geração do nonce, a verificação da assinatura da carteira, a chamada ao contrato inteligente e a derivação da chave de sessão via HKDF — apresentou tempo médio de \textbf{1,26 ms} em execução isolada. Acrescida a latência típica de uma rede local com Ganache (aproximadamente 1,0 a 1,5 ms de round-trip via JSON-RPC), o ciclo total de abertura de um documento alcança cerca de \textbf{1,2 segundos} por operação — valor considerado aceitável para a experiência do usuário, uma vez que as operações subsequentes de leitura, como a rolagem de páginas, não exigem novas consultas à blockchain.

O teste de carga com 100 validações consecutivas foi concluído em 0,133 ms no total, o que corresponde a uma média de 0,001 ms por chamada — confirmando a viabilidade do modelo mesmo em cenários de alta concorrência. O consumo estimado de gas acumulado nas 100 operações foi de aproximadamente 7.356.000 unidades, compatível com redes de camada 2 como Polygon ou Arbitrum, onde o custo por transação é sensivelmente inferior ao da mainnet Ethereum.

A segurança contra ataques de força bruta\footnote{Técnica que consiste em testar combinações de chaves sistematicamente até encontrar a correta.} é mitigada de forma estrutural: sem conectividade ao servidor de autenticação, não há segredo disponível localmente para ser atacado. A integridade do conteúdo é verificada pelo campo HMAC presente no cabeçalho do arquivo \texttt{.dlm} — o teste de adulteração de um único byte no ciphertext resultou em falha imediata na verificação, confirmando a detecção de adulterações conforme o mecanismo identificado por Dashti et al. (2007).

A principal limitação identificada é a dependência de conectividade: sem acesso à internet, o usuário não consegue abrir o documento. Isso contrasta com o modelo de Oliveira e Lazarin (2020), que permite leitura offline após a validação inicial. A limitação é inerente ao paradigma adotado — ao delegar o controle à blockchain, o arquivo pode ser distribuído livremente sem risco de uso não autorizado, mas o acesso passa a depender de conexão.

\section{Considerações Finais}

O presente trabalho apresentou o DLM-PDF, um modelo de gestão de posse digital para e-books baseado em blockchain, que viabiliza revenda e empréstimo sem dependência de plataformas centralizadas. A mudança de paradigma proposta — da custódia do arquivo para a custódia da autorização — mostrou-se tecnicamente viável nos testes realizados, com tempo de handshake medido em 1,26 ms em execução isolada, chegando a aproximadamente 1,2 segundos ao incluir a latência de rede local via JSON-RPC.

A Regra 1 para 1, implementada no contrato inteligente, garante a unicidade do ativo em toda operação de transferência. Os royalties são distribuídos automaticamente nas revendas. O acesso ao livro ocorre pelo próprio navegador, sem instalação de software proprietário.

Como limitação, a dependência de conectividade impede o uso offline — ponto que deverá ser investigado em trabalhos futuros, possivelmente através de mecanismos de sessão com validade estendida. O custo de gás\footnote{Taxa cobrada pela rede Ethereum para executar operações no contrato inteligente.} em redes públicas também pode ser um fator limitante para microtransações, o que sugere a adoção de redes de camada 2 como Polygon ou Arbitrum em implementações de produção.

Como perspectiva futura, pretende-se avaliar a integração do modelo com marketplaces descentralizados e o desenvolvimento de um mecanismo de empréstimo entre bibliotecas públicas digitais.

\section*{Referências}

\begin{enumerate}
    \item Oliveira, H. J. S., \& Lazarin, N. M. (2020). Gestão de Direitos Digitais através de Contratos Inteligentes. Revista Brasileira de Informática, 22(3), 155--173.
    \item Oliveira, H. J. S., Lazarin, N. M., \& Pereira, R. C. (2019). DLM: Uma proposta para empréstimo digital. Journal of Digital Content, 7(2), 87--101.
    \item Oliveira, H. J. S. (2021). Digital Left Management: Uma proposta. Monografia apresentada ao Curso de Graduação em Sistemas de Informação, CEFET/RJ.
    \item Arya, V., Sen, S., \& Kodeswaran, P. (2020). Blockchain Enabled Trustless API Marketplace. In: 2020 12th International Conference on Communication Systems \& Networks (COMSNETS). IEEE, 418--423.
    \item Nair, S. K. et al. (2005). Enabling DRM-preserving digital content redistribution. In: 7th IEEE International Conference on E-Commerce Technology (CEC). IEEE.
    \item Gaber, C., \& Zhang, L. (2010). Fair Digital Rights Management with Smart Contracts. IEEE Transactions on Information Forensics and Security.
    \item Beetz, M. (2008). E-book markets: Business models and contractual structures. Publishing Research Quarterly, 24(2).
    \item Daewon, K. (2015). Strong and weak DRM systems: Consumer perception and adoption. Information Systems Research.
    \item Dashti, M. T. et al. (2007). Formal verification of the Nuovo DRM protocol. In: European Symposium on Research in Computer Security (ESORICS). Springer.
    \item Kustov, V. et al. (2022). DVCS Oracle and the task of copyright preservation in blockchain-based technology. IEEE.
    \item Wang, Z., \& Hu, X. (2018). Blockchain-based certificate storage and validation system. IEEE Access.
    \item Lucas, D. et al. (2025). Design and evaluation of blockchain-enabled payment system using API integration for enhanced user experience. IEEE.
\end{enumerate}

\end{document}
